\chapter*{函数综合问题练习}
\begin{enumerate}
    \item 整根问题
    \begin{enumerate}
        \item 若关于$x$的不等式$k|x| > |x - 2|$恰有4个整数解，则实数$k$的取值范围是\blkda{$(\frac{3}{5}, \frac{2}{3}]$}
        \begin{jieda}
            易知$x = 0$不是不等式的解，故$k|x| > |x - 2|$进行参变分离，同解变形$k > |\frac{x-2}{x}|$，即$k > |1 - \frac{2}{x}|$ \\
            由函数$f(x) = |1 - \frac{2}{x}|$的图像可知整根的集合为$\{2, 3, 4, 5\}$，故$f(5) < k \leq f(6)$，即$k \in (\frac{3}{5}, \frac{2}{3}]$
            \ps \textbf{解决整根问题的第一步往往是确认整根的集合}
        \end{jieda}
        \item 已知关于$x$的不等式组$\left \{ \begin{array}{l}
             \frac{x-a}{x-1+a} < 0\\
             x+3a>1 
        \end{array}\right.$的整数解恰有两个，则实数$a$的取值范围是\blkda{$(1, 2]$}
        \begin{jieda}
            首先可以确定第一个不等式的解集为$(a, 1-a)$或$(1-a, a)$，因为$a$和$1-a$关于$\frac{1}{2}$对称，因此$\frac{1}{2}$必然在该不等式的解集中。而第二个不等式的解集为$(1-3a, +\infty)$
            \begin{dingautolist}{172}
                \item 若$a \in [0, 1]$，则$(a, 1-a)$或$(1-a, a)$没有整数解，故整个不等式组的解集中也没有整数解.
                \item 若$a < 0$，则$a < 0 < 1 < 1-a < 1-3a$，则整个不等式组解集为空集.
                \item 若$a > 1$，则$1-3a < 1-a < 0 < 1 < a$整个不等式组的解集为$(1-a, a)$，此时解集中必有两个整数解$0, 1$，故当$a \in (1, 2]$时，满足要求.
            \end{dingautolist}
            综上，$a \in (1, 2]$ \\
            \ps 本题画数轴可以更方便分析，本题判断出$a$和$1-a$关于$\frac{1}{2}$对称后可以猜测到两个整根是0和1.
        \end{jieda}
        \item 已知不等式组$ \left\{ \begin{array}{l}
             x^2-x+a-a^2 < 0  \\
             x + 2a > 1
        \end{array} \right.$的整数解恰好有两个，$a$的取值范围是\blkda{$(1, 2]$}
        \begin{jieda}
            设函数$f(x) = x^2 - x$，则第一个不等式即$f(x) < f(a)$，根据函数的性质可知不等式即$|x-\frac{1}{2}| < |a - \frac{1}{2}|$，故$\frac{1}{2}-|a - \frac{1}{2}| < x < \frac{1}{2} + |a - \frac{1}{2}|$，故第一个不等式解集为$(1 - a, a)$或$(a, 1-a)$，第二个不等式解集为$(1-2a, +\infty)$
            \begin{dingautolist}{172}
                \item $a \in [0, 1]$，此时第一个不等式中不含有两个整数解，不符合要求.
                \item $a < 0$，第一个不等式解集为$(a, 1-a)$，而$1-a < 1-2a$，故不等式组解集为空集，不符合要求.
                \item $a > 1$，第一个不等式解集为$(1 - a, a)$，故不等式解集为$(1 - a, a)$，当$a \in (1, 2]$满足要求.
            \end{dingautolist}
        \end{jieda}
    \end{enumerate}
    \item 变量联动与不联动的区别
    \begin{enumerate}
        \item 已知函数$f(x) = x - 2$，$g(x) = mx^2-2mx+1(m \in \mathbf{R})$. 
        \begin{enumerate}
            \item 对任意的$x \in \mathbf{R}$，不等式$g(x) > f(x)$恒成立，求$m$取值范围.
            \item 对任意的$x_1 \in [1, 2], x_2 \in [-1, 0]$，不等式$g(x_1) > f(x_2)$恒成立，求$m$取值范围.
            \item 对任意的$x_1 \in [1, 2]$，存在$x_2 \in [3, 4]$，使得$g(x_1) = f(x_2)$，求$m$取值范围.
            \item 对任意的$x_1 \in [1, 2]$，存在$x_2 \in [0, 2]$，使得$g(x_1) > f(x_2)$，求$m$取值范围.
        \end{enumerate}
        
        \begin{jieda}
            \begin{enumerate}
                \item 问题转化为$g(x) - f(x) > 0$在$\mathbf{R}$上恒成立，即$mx^2-(2m+1)x+3 > 0$恒成立. $m = 0$不满足要求，$m < 0$不满足要求，$m > 0$时应有$\Delta < 0$，即$m \in (\frac{2-\sqrt{3}}{2}, \frac{2+\sqrt{3}}{2})$
                \item 问题转化为$g(x)$在$[1, 2]$上的最小值大于$f(x)$在$[3, 4]$上的最大值. $f(x)$在$[-1, 0]$上的最大值为$-2$，故问题转化为$g(x) + 2 > 0$在$[1, 2]$上恒成立，设$h(x) = g(x) + 2 = mx^2-2mx + 3 = m(x-1)^2-(m-3)$，下面分析$h(x)$的最值
                \begin{dingautolist}{172}
                    \item $m > 0$，$h(x)$在$[1, 2]$内单调递增，因此最小值为$h(1) = 3-m$
                    \item $m = 0$，$h(x) = 3$，最小值为$3$
                    \item $m < 0$，$h(x)$在$[1, 2]$内单调递减，因此最小值为$h(2) = 3$
                \end{dingautolist}
                故当$m \in (-\infty, 3)$时，满足要求.
                \item 问题转化为$g(x)$在$[1, 2]$上的值域是$f(x)$在$[3, 4]$上值域的子集. $f(x)$在$[3, 4]$上值域为$[1, 2]$，$g(x) = m(x-1)^2 + (1-m)$，下面分析$g(x)$的值域
                \begin{dingautolist}{172}
                    \item $m > 0$，$g(x)$在$[1, 2]$内单调递增，值域为$[1-m, 1]$
                    \item $m = 0$，$g(x) = 1$，值域为$\{1\}$
                    \item $m < 0$，$g(x)$在$[1, 2]$内单调递增，值域为$[1, 1-m]$
                \end{dingautolist}
                故当$m \in [-1, 0]$时满足要求.
                \item 问题转化为$g(x)$在$[1, 2]$上的最小值大于$f(x)$在$[0, 2]$上的最小值. $f(x)$在$[0, 2]$上最小值为-2，下面分析$g(x)$的最值
                \begin{dingautolist}{172}
                    \item $m > 0$，$g(x)$在$[1, 2]$内单调递增，因此最小值为$g(1) = 1-m$
                    \item $m = 0$，$g(x) = -1$，最小值为$-1$
                    \item $m < 0$，$g(x)$在$[1, 2]$内单调递减，因此最小值为$g(2) = 1$
                \end{dingautolist}
                故当$m \in (-\infty, 3)$时，满足要求.
            \end{enumerate}
        \end{jieda}
        \item 已知函数 $f\left( x\right)  = {\left( \dfrac{1}{2}\right) }^{x}$ 与 $g\left( x\right)  = {\log }_{4}\left( {{x}^{2} - {2ax} + 4}\right) \left( {a > 0}\right)$ ,若对任意的 ${x}_{1} \in  \left( {0,1}\right)$ ,都存在 ${x}_{2} \in  \left\lbrack  {0,2}\right\rbrack$ ,使得 $f\left( {x}_{1}\right)  = g\left( {x}_{2}\right)$ ,则实数 $a$ 的取值范围是\blkda{$[\sqrt{2}, +\infty)$}
        \begin{jieda}
            问题转化为$f(x)$在$(0, 1)$上的值域是$g(x)$在$[0, 2]$上值域的子集. \\
            $f(x)$在$(0, 1)$上的值域为$\left(\dfrac{1}{2}, 1\right)$，因此内层函数$h(x) = {x}^{2} - {2ax} + 4$在$[0, 2]$上的值域要包含$(2, 4)$，因此内层函数对称轴为直线$x = a$
            \begin{dingautolist}{172}
                \item $a \in (0, 2]$ \\
                此时$h(x)$在$[0, a]$上单调减，在$[a, 2]$上单调增，因此当$h(a) \leq 2$时满足要求，即$4-a^2\leq2$，解得$a \in [\sqrt{2}, 2]$
                \item $a > 2$ \\
                此时$h(x)$在$[0, 2]$上单调减，因此当$h(2) \leq 2$时满足要求，即$0 < 8-4a \leq 2$，在$a > 2$时恒成立.
            \end{dingautolist}
            综上，$a \in [\sqrt{2}, +\infty)$
        \end{jieda}
    \end{enumerate}
\item 识图
\begin{enumerate}
    \item 函数$y=\dfrac{x-2}{x-1}$的图像是\dotfill(\kh{B})
    \fourtu
    {\includegraphics[totalheight=4.5\baselineskip]{./bodyxb/src/bx01ch03hanshudetuxiangt8a.png}}
    {\includegraphics[totalheight=4.5\baselineskip]{./bodyxb/src/bx01ch03hanshudetuxiangt8b.png}}
    {\includegraphics[totalheight=4.5\baselineskip]{./bodyxb/src/bx01ch03hanshudetuxiangt8c.png}}
    {\includegraphics[totalheight=4.5\baselineskip]{./bodyxb/src/bx01ch03hanshudetuxiangt8d.png}}
    \begin{jieda}
        可以先确定对称中心是$(1, 1)$，然后再代入$x = 0$判断.
    \end{jieda}

    \item 函数$y = (|x|+1)ln|x|$的图象大致为\dotfill{(\kh{B})}
    \begin{figure}[h]
        \centering
        \includegraphics[width=0.5\linewidth]{q34.png}
    \end{figure}
    \begin{jieda}
        首先可以判断出来该函数为偶函数，其次可以判断出来在$x > 0$时函数是单调递增的，因此选B
    \end{jieda}
    \item 函数$y = \dfrac{x^2}{2^x-2^{-x}}$的图象大致为\dotfill{(\kh{A})} \\
    \begin{figure}[h]
        \centering
        \includegraphics[width=0.5\linewidth]{q35.png}
    \end{figure} \\
    \begin{jieda}
        首先可以判断出来该函数是奇函数，并且当$x > 0$时$y >0$。通过带入特值($x=1, x=\dfrac{1}{2}$)（或求导）等可以判断出来函数在$x > 0$时不是单调递减的，因此选A
    \end{jieda}
    \item 已知函数$y=f(\abs{x})$的图像如左下图所示, 则函数$y=f(x)$的图像不可能是\dotfill(\kh{B})
    \begin{minipage}{\linewidth}
    \centering
    \includegraphics[height=5.0\baselineskip]{./bodyxb/src/bx01ch03hanshudetuxiangt9start.png}\qquad
    
\includegraphics[height=5.0\baselineskip]{./bodyxb/src/bx01ch03hanshudetuxiangt9a.png}\qquad
    
\includegraphics[height=5.0\baselineskip]{./bodyxb/src/bx01ch03hanshudetuxiangt9b.png}\qquad
    
\includegraphics[height=5.0\baselineskip]{./bodyxb/src/bx01ch03hanshudetuxiangt9c.png}\qquad
    
\includegraphics[height=5.0\baselineskip]{./bodyxb/src/bx01ch03hanshudetuxiangt9d.png}
    \end{minipage}
    \begin{jieda}
        $f(|x|)$是将$f(x)$在$y$轴右侧图像翻到左侧
    \end{jieda}
    
    \item 已知函数$f(x)=ax^3+bx^2+cx+d$的图像如下图, 则\dotfill(\kh{A})
    \xx{$b\in(-\infty,0)$}{$b\in(0,1)$}{$b\in(1,2)$}{$b\in(2,+\infty)$}
    \begin{minipage}{\linewidth}
    \centering
    \includegraphics[height=5\baselineskip]{./bodyxb/src/bx01ch03hanshudetuxiangt11.png}
    \end{minipage}
    \begin{jieda}
        首先可以判断出$a > 0$，然后根据三个零点列方程可以得到$b = -3a$或者将$f(x)$写成因式相乘的形式：$f(x) = ax(x-1)(x-2)$，也可以得到$b = -3a$
    \end{jieda}
    \item 任取$x_1,x_2\in[a,b], x_1\neq x_2$, 若$f\left(\dfrac{x_1+x_2}{2}\right)>\dfrac{f(x_1)+f(x_2)}{2}$, 则称$f(x)$是$[a,b]$上的凸函数. 
    下列图像中, 可以是凸函数图像的是\dotfill(\kh{C})
    \fourtu
    {
\includegraphics[totalheight=5.5\baselineskip]{./bodyxb/src/bx01ch03hanshudetuxiangt10a.png}}
    {
\includegraphics[totalheight=5.5\baselineskip]{./bodyxb/src/bx01ch03hanshudetuxiangt10b.png}}
    {
\includegraphics[totalheight=5.5\baselineskip]{./bodyxb/src/bx01ch03hanshudetuxiangt10c.png}}
    {
\includegraphics[totalheight=5.5\baselineskip]{./bodyxb/src/bx01ch03hanshudetuxiangt10d.png}}
    \item 已知函数$y=f(x)$与$y=g(x)$在$[-2,2]$的图像如图所示: \\
    \begin{minipage}{\linewidth}
    \centering
    
\includegraphics[totalheight=6\baselineskip]{./bodyxb/src/bx01ch03hanshudeyunsuan02t1401.png}
    \qquad\qquad
    
\includegraphics[totalheight=6\baselineskip]{./bodyxb/src/bx01ch03hanshudeyunsuan02t1402.png}
    \end{minipage}
    
    判断下列命题的真假, 并说明理由: 
    \begin{enumerate}
    \item 方程$f(g(x))=0$有且仅有6个根; 
    \item 方程$g(f(x))=0$有且仅有3个根; 
    \item 方程$f(f(x))=0$有且仅有5个根; 
    \item 方程$g(g(x))=0$有且仅有4个根. 
    \end{enumerate}
    \begin{jieda}
    $f(x)=0$的根可取近似值为$0,\pm1.5$; $g(x)=0$的根可取近似值为$-1.8,0.4$. 
    \begin{enumerate}
    \item $f(g(x))=0\iff g(x)=0,\pm1.5$, 结合$g(x)$的图像可得解的个数为$3+3=6$个. 
    \item $g(f(x))=0\iff f(x)=-1.8,0.4$, 结合$f(x)$的图像可得解的个数为$1+3=4$个. 
    \item $f(f(x))=0\iff f(x)=0,\pm1.5$, 结合$f(x)$的图像可得解的个数为$3+2=5$个. 
    \item $g(g(x))=0\iff g(x)=-1.8,0.4$, 结合$g(x)$的图像可得解的个数为$2+2=4$个. 
    \end{enumerate}
    综上, 只有(2)是假命题, 其他都是真命题. 
    \end{jieda}
    \end{enumerate}
    \item 函数的零点、交点与方程的根
    \begin{enumerate}
        \item 用$C(A)$表示非空集合A中元素个数，定义$|A-B| = \left \{ \begin{array}{ll}
             C(A)-C(B),& C(A) \geq C(B)  \\
             C(B)-C(A),& C(A) < C(B)
        \end{array}\right.$，
        已知集合$A = \{x | x^2 + x = 0\}, B = \{x | (3x^2+ax)(x^2+ax+2) = 0\}$，且$|A-B|=1$，设实数$a$的所有可能取值构成集合$S$，则$C(S) = $\blkda{5}
        \begin{jieda}
            $A = \{0, -1\}, C(A) = 2$,因为$|A-B| = 1$，所以$C(B) = 1 \mbox{或} 3$。 \\
            方程$(3x^2+ax)(x^2+ax+2) = 0$的根包括$0, -\dfrac{a}{3}, \dfrac{-a+\sqrt{a^2-8}}{2}, \dfrac{-a-\sqrt{a^2-8}}{2}$ 
            \begin{enumerate}
                \item 当$a \in (-2\sqrt{2}, 2\sqrt{2})$时，后两者并非实根，因此$C(B) \neq 3$，当$a = 0$的时候$B = \{0\}, C(B) = 1$
                \item 当$a \in \{-2\sqrt{2}, 2\sqrt{2}\}$时，$B = \{0, -\dfrac{a}{3}, -\dfrac{a}{2}\}, C(B) = 3$
                \item 当$a \in (-\infty, -2\sqrt{2}) \cup (2\sqrt{2}, +\infty)$时，若有$-\dfrac{a}{3} = \dfrac{-a \pm \sqrt{a^2-8}}{2}$可使得$C(B) = 3$，解得$a = \pm 3$
            \end{enumerate}
            综上$S = \{0, 2\sqrt{2}, -2\sqrt{2}, 3, -3\}$，因此$C(S) = 5$
        \end{jieda} 
        \item 用$C(A)$表示非空集合$A$中的元素个数，定义$|A-B| = \left \{ \begin{array}{ll}
            C(A) - C(B), & C(A) \geq C(B) \\
            C(B) - C(A), & C(A) < C(B)
        \end{array}\right.$，若$A = \{1,2\}, B = \{x | |x^2+2x-3| = a\}$，且$|A-B| = 1$，由$a$的所有可能值构成的集合为$S$，那么$C(S) = $\blkda{1}
        \begin{jieda}
            $|x^2+2x-3| = a$即$y = a$与$y = |x^2+2x-3|$的交点横坐标，而$|A-B| = 1$即两者交点数为3或1. 易知两者交点个数为0,2,3,4，故考虑其有三个交点的情况，即$a = 4$，故$S$中只有一个元素.
        \end{jieda}
        \item 已知方程$|x| = ax+1$有一负解，且无正解，则实数$a$的取值范围是\blkda{$[1, +\infty)$ }
        \begin{jieda}
            将方程的根转化为$y = |x|$与$y = ax+1$的交点的横坐标.$y = ax+1$过$(0, 1)$，斜率为$a$，故当$a \in [1, +\infty)$时，两函数在$x > 0$时无交点，而在$x < 0$时有一个交点，满足要求.
        \end{jieda}
        \item 设$f(x)$是定义在$\bR$上的奇函数, 且对于任意的$x\zr$, $f(1+x)-f(1-x)=0$恒成立; 当$x\in[0,1]$时, $f(x)=2x$. 
        若方程$f(x)=ax$恰好有5个不同的解, 则实数$a$的取值范围为\blkda[2.5]{$(-\frac{2}{3},-\frac{2}{7})\cup\{\frac{2}{5}\}$}. 
        \begin{jieda}
        $f(x)$图像关于原点对称, 关于直线$x=1$对称, 可知周期$T=4$. \\
        $y = ax$的图像是过原点的直线，随着$a$的变化，其斜率变化. 绘制草图后可以判断要分为$a>0$与$a<0$两种情形, 
        不难求得$a=\frac{2}{5}$或$a\in(-\frac{2}{3},-\frac{2}{7})$. 
        \end{jieda}
        \item 已知$f(x) = x^2 + \frac{1}{x^2} + k (x + \frac{1}{x})+11, x \in (0, +\infty)$，若$f(x) = 0$有两个解，则实数$k$的取值范围是\blkda[4]{$(-\infty, -\frac{13}{2}) \cup \{-6\}$}
        \begin{jieda}
            首先考虑换元$t = x + \frac{1}{x}$，$f(x) = g(t) = t^2+k\cdot t + 9$，当$t = 2$时，只对应一个$x$，当$t > 2$时，对应着两个$x$. 故若$f(x)$在$(0, +\infty)$，有两个零点可以转化为$g(t)$在$(2, +\infty)$有一个零点，进一步转化为$y = -k$与$y = t + \frac{9}{t}$在$(2, +\infty)$有一个交点. 因此$k \in (-\infty, -\frac{13}{2}) \cup \{-6\}$ \\
            \ps \textbf{换元后及时判断新元的取值范围，并且分析清楚换元前后的对应关系！}
        \end{jieda}
        \item 若关于$x$的方程$\frac{|x|}{x+6} = kx^2$有四个不同的实数解，求实数$k$的取值范围.
        \begin{jieda}
            易知必有一根为0 \\
            参变分离有$\frac{1}{k} = \frac{(x+6)x^2}{|x|} = \left \{ \begin{array}{ll}
                \frac{(x+6)x^2}{x}, & x > 0 \\
                -\frac{(x+6)x^2}{x},  & x < 0
            \end{array}\right.$
            由于原方程必有一根为0，故可以转化为方程$\frac{1}{k} = \left \{ \begin{array}{ll}
                (x+6)x, & x > 0 \\
                -(x+6)x,  & x < 0
            \end{array}\right.$有三个不等实根.因此$\frac{1}{k} \in (0, 9)$，即$k \in (\frac{1}{9}, +\infty)$
        \end{jieda}
        \item 关于$x$的方程$k\cdot9^x-k\cdot3^{x+1}+6(k-5)=0$在$[0,2]$上有唯一解, 求实数$k$的取值范围. 
        \begin{jieda}
        令$t=3^x\in[1,9]$, $t$与$x$一一对应.于是关于$t$的方程$kt^2-3kt+6(k-5)=0$在$[1,9]$上有唯一解. 于是$\frac{30}{k}=t^2-3t+6$. 结合图像可得$\frac{30}{k}\in(4,60]\cup\{\frac{15}{4}\}$, 即$k\in[\frac{1}{2},\frac{15}{2})\cup\{8\}$. 
        \end{jieda}
        \item 已知关于$x$的方程$(2^x-1)^2-2\cdot\abs{2^x-1}+k=0$有三个根, 实数$k$的取值范围是\blkda{$(0, 1)$}. 
        \begin{jieda}
            换元$t = 2^x - 1$，$t \in (-1, +\infty)$，$t$与$x$一一对应. 故原方程有三个解即$t^2 - 2|t| + k = 0$有三个根，即$y = k$与$y = 2|t| - t^2$有三个交点.函数$y = 2|t| - t^2$在$(-1, 0)$单调递减，值域为$(0, 1)$，在$(0, 1)$单调递增，值域为$(0, 1)$，在$[1, +\infty)$单调递减，值域为$(-\infty, 1]$，故$k \in (0, 1)$时两函数有三个交点.
        \end{jieda}
    \item 对于实数$a, b$定义运算“*”：$a*b = \left\{  
         \begin{array}{ll}  
            a^2 - ab, & a \leq b, \\  
            b^2 - ab, & a > b.
         \end{array}
    \right.$ 设$f(x) = (2x-1) * (x-1)$，且关于$x$的方程$f(x) = m (m \in \mathbf{R})$恰有三个互不相等的实数根$x_1, x_2, x_3$，则$x_1x_2x_3$的取值范围是\blkda[3]{$(-\frac{\sqrt{3} - 1}{16}, 0)$}
    \begin{jieda}
        根据运算定义规则，首先求$2x-1$和$x-1$的交点，为$(0, -1)$，\\
        故$f(x) = \left\{ 
                 \begin{array}{ll}  
                    (2x-1)^2 - (2x-1)(x-1), & x \leq 0, \\  
                    (x-1)^2 - (2x-1)(x-1), & x > 0.
                 \end{array}
                \right. $ \\
        即$f(x) = \left\{ 
                 \begin{array}{ll}  
                    2x^2-x, & x \leq 0, \\  
                    x-x^2, & x > 0.
                 \end{array}
                \right. $ \\
        观察函数图像可知，在满足题意的情况下，$y = m (0 < m < \frac{1}{4})$与$f(x)$在$y$轴左侧有一个交点，在$y$轴右侧有两个交点。故由韦达定理有$x_2x_3 = m$，$x_1 = \frac{1-\sqrt{1+8m}}{4}$。故$x_1x_2x_3 = m \cdot \frac{1-\sqrt{1+8m}}{4} = -\frac{1}{4}(m \cdot (\sqrt{1+8m} - 1))$。因为$-\frac{1}{4}(m \cdot (\sqrt{1+8m} - 1))$在$m \in (0, \frac{1}{4})$单调递减，故其取值范围是$(-\frac{\sqrt{3} - 1}{16}, 0)$
    \end{jieda}
    \item 已知函数$f(x) = \left\{  
         \begin{array}{ll}  
            x^3, & x \geq 0, \\  
            -x, & x < 0.
         \end{array}
    \right.$ 若函数$g(x) = f(x) - |kx^2-2x| (k \in \mathbf{R})$恰有4零点，$k$的取值范围是\blkda[3.5]{$(-\infty, 0) \cup (2\sqrt{2}, +\infty)$}.
    \begin{jieda}
        容易判断，$g(x)$有一零点为0. 设$F(x) = \left\{  
         \begin{array}{ll}  
            x^2, & x \geq 0, \\  
            1, & x < 0.
         \end{array}
        \right.$，$H(x) = |kx-2|$，故其余三个零点是函数$G(x) = F(x) - H(x)$的零点.易知$k = 0$时$G(x)$没有三个零点. 
        \begin{enumerate}
            \item 当$k < 0$时，$G(x)$在$(-\infty, \frac{2}{k})$单调增，在$(\frac{2}{k}, 0)$单调减，在$(0, -\frac{k}{2})$单调减，在$(-\frac{k}{2}, +\infty)$单调增.$G(\frac{3}{k}) = 0, G(\frac{1}{k}) = 0, G(0) = -2, G(2-k) = 2-2k$，故$G(x)$有三个零点，分别是$\frac{3}{k}, \frac{1}{k}$，以及一个位于$(0, 2-k)$的零点.
            \item 当$k > 0$时，当$x < 0$，$G(x) = kx-1$，无零点. \\
            当$x \in (0, \frac{2}{k})$，$G(x)$单调增,$G(\frac{1}{k}) = 0$. \\
            当$x \geq \frac{2}{k}$，$G(x) = x^2 - kx + 2$，要使得$G(x)$在$[\frac{2}{k}, +\infty)$有两个零点，必有$\Delta > 0$，即$k > 2\sqrt{2}$. 抛物线的对称轴为$x = \frac{k}{2}$，因为当$k > 2\sqrt{2}$时$\frac{k}{2} > \frac{2}{k}$，又因为$G(\frac{2}{k}) > 0$，所以当$k > 2\sqrt{2}$时，$G(x)$在$[\frac{2}{k}, +\infty)$有两个零点.
        \end{enumerate}
        综上，$k \in (-\infty, 0) \cup (2\sqrt{2}, +\infty)$ \\
        \ps 从数形结合角度平移$H(x)$来看会更直观.
    \end{jieda}
    \item 已知关于$x$的方程$(x^2-1)^2-|x^2-1|+k = 0$. 下列语句中说法正确的有\blkda{\ding{172} \ding{173} \ding{174} \ding{175}}
    \begin{dingautolist}{172}
        \item 存在实数$k$使得方程恰有2个不同的实数解.
        \item 存在实数$k$使得方程恰有4个不同的实数解.
        \item 存在实数$k$使得方程恰有5个不同的实数解.
        \item 存在实数$k$使得方程恰有8个不同的实数解.
    \end{dingautolist}
    \begin{jieda}
        换元$t = |x^2-1|$，参变分离有$k = t - t^2$
        对于内层方程$t = |x^2-1|$：
        \begin{enumerate}
            \item 若$t < 0$，则方程无解
            \item 若$t = 0$，则方程有2个解
            \item 若$t \in (0, 1)$，则方程有4个解
            \item 若$t = 1$，则方程有3个解
            \item 若$t > 1$，则方程有2个解
        \end{enumerate}
        对于外层方程$k = t - t^2$：
        \begin{enumerate}
            \item 若$k > \dfrac{1}{4}$，则外层方程无解，复合方程无解
            \item 若$k = \dfrac{1}{4}$，则存在唯一的$t = \dfrac{1}{2}$使外层方程成立，对复合方程本身有4个解
            \item 若$0 < k < \dfrac{1}{4}$，则外层方程两解$t_1, t_2 \in (0, 1)$，复合方程有8个解
            \item 若$k = 0$，则外层方程两解$t_1 = 0, t_2 = 1$，复合方程有5个解
            \item 若$k < 0$，则外层方程两解$t_1 < 0, t_2 > 1$，复合方程有2个解
        \end{enumerate}
    \end{jieda}
    \end{enumerate}
    \item 利用必要条件简化问题
    \begin{enumerate}
        \item 已知函数$y=\frac{1}{4-2^x}$的图像关于点$P$对称, 则点$P$的坐标为\blkda{$(2,\frac{1}{8})$}
        \begin{jieda}
        设$P(a,b)$, 因为$x\neq 2$, 故$a=2$. \\
        $2b=f(1)+f(3)$, 于是$b=\frac{1}{8}$. \\
        \ps $f(x)$的对称中心是$(a, b)$或对称轴是$x = a$的必要条件均为$f(x)$的定义域关于$a$对称.
        \end{jieda}
        \item 已知二次函数 $f(x)=-\frac{1}{2} x^2+x$，是否存在实数 $m, n(m<n)$, 使得 $f(x)$ 的定义域为 $[m, n]$, 值域为 $[k m, k n](k>1)$.
        \begin{jieda}
        $f(x)=-\dfrac{1}{2} x^2+x=-\dfrac{1}{2} \times(x-1)^2+\dfrac{1}{2} \leq \dfrac{1}{2}$
        $\therefore k n \leq \dfrac{1}{2}, \therefore m<n<\dfrac{1}{2k}$
        又 $\because$ 函数的对称轴为 $x=1$
        $\therefore f(x)$ 在 $[m, n]$ 单调递增则有 $f(m)=k m, f(n)=k n$
        故$m, n$是$-\frac{1}{2}x^2 + x = kx$的两根，故$m = 2(1-k), n = 0$ \\
        \ps 根据值域判断函数在定义域内的单调性是简化，解决问题的突破口。\\
        常规解决方法是分析单调性，根据最值列方程组.
        \end{jieda}
        \item 若函数$f(x)$在区间$[a, b]$上的值域恰是$\left[\dfrac{1}{b}, \dfrac{1}{a}\right]$，则称区间$[a, b]$是$f(x)$的一个倒域区间. 求函数$g(x) = \left \{ \begin{array}{ll}
            -x^2+2x & 0 \leq x \leq 2 \\
            x^2 + 2x & -2 \leq x < 0
        \end{array}\right.$的所有倒域区间.
        \begin{jieda}
            $g(x)$的定义域为$[-2, 2]$，$g(x)$值域为$[-1, 1]$，故有$-2 \leq a < b \leq 2, -1 \leq \frac{1}{b} < \frac{1}{a} \leq 1$，解得$1 \leq a < b \leq 2$或$-2 \leq a < b \leq -1$.
            \begin{dingautolist}{172}
                \item  $1 \leq a < b \leq 2$，此时$g(x)$单调递减，有$\left \{ \begin{array}{l}
                     -a^2 + 2a = \frac{1}{a}  \\
                     -b^2 + 2b = \frac{1}{b} 
                \end{array} \right.$，即$a, b$是$-x^2 + 2x = \dfrac{1}{x}$的两根，即$a, b$是$-x^3 + 2x^2 -1 = 0$的两根，可以看出有一根为$1$，即$a, b$是$(x-1)(-x^2+x+1) = 0$的两根，故$a = 1, b = \dfrac{1+\sqrt{5}}{2}$
                \item $-2 \leq a < b \leq -1$，此种情况和前一种对称，可得$a = -\dfrac{1+\sqrt{5}}{2}, b = -1$
            \end{dingautolist}
            因此两个倒域区间为$\left[1, \dfrac{1+\sqrt{5}}{2}\right]$和$\left[-\dfrac{1+\sqrt{5}}{2}, -1\right]$
        \end{jieda}
        \item 若函数$f(x)$在区间$[m, n]$上的值域恰是$[m, n]$，则称区间$[m, n]$是$f(x)$的一个等域区间. 若区间$[m,n]$是函数$y = \dfrac{(a+1)x - \dfrac{1}{a}}{ax}(a \neq 0)$的一个等域区间.求$n-m$的最大值.
        \begin{jieda}
            $\dfrac{(a+1)x - \dfrac{1}{a}}{ax} = \dfrac{a+1}{a} - \dfrac{1}{a^2x}$，故函数为$y = \dfrac{a+1}{a} - \dfrac{1}{a^2x} (a \neq 0)$，易知函数的定义域为$(-\infty, 0) \cup (0, +\infty)$，故必有$m < n < 0$或$0 < m < n$. 故易知函数在$[m, n]$内严格单调递增，故$\left \{ \begin{array}{l}
                    \dfrac{a+1}{a} - \dfrac{1}{a^2m} = m \\
                    \dfrac{a+1}{a} - \dfrac{1}{a^2n} = n
                \end{array} \right.$，即$m, n$是$a^2x^2 - (a+1)ax + 1 = 0$的两根. 由韦达定理有$\left \{ \begin{array}{l}
                    m+n = \dfrac{a+1}{a}\\
                    m \cdot n = \dfrac{1}{a^2}
                \end{array} \right.$，由$\Delta > 0$得$((a+1)^2 - 4)a^2 > 0$，故$a \in (-\infty, -3]\cup [1, +\infty)$ \\
                $n - m = \sqrt{(m+n)^2 - 4mn} = \sqrt{(\dfrac{a+1}{a})^2 - \dfrac{4}{a^2}} = \sqrt{\dfrac{a^2+2a-3}{a^2}} \leq \dfrac{2\sqrt{3}}{3}$，当$a = 3$时取等号.
        \end{jieda}
    \end{enumerate}
\end{enumerate}